% section: 式の変形

  \section{式の変形}
    \subsection{$+a-a$をする}
      平方完成を例に考えよう.
      平方完成とは,ある二次式
      \[
        ax^2+bx+c
      \]
      に対して
      \begin{align*}
        ax^2+bx+c&=a\left(x^2+\frac{b}{a}x\right) +c\\
        &=a\left(x^2+\frac{b}{a}x+\frac{b^2}{4a^2}-\frac{b^2}{4a^2}\right)\\
        &\qquad +c \tag*{\(\cdots\) (☆)}\\
        &=a\left(x+\frac{b}{2a}\right)^2-\frac{b^2}{4a}+c
      \end{align*}
      とすることで$x$に関する項を1つにまとめる操作のことである.
      ここで注目してほしいのは☆部分の式変形である.
      これが意味するのは,無理やりに2乗の形,すなわち平方を作り上げよう,ということだ.
      これにより$x$に関する項は定数とともに括弧の中に閉じ込められ,よりシンプルに式が記述できる.
      二次関数の最大・最小問題や軸の決定などにおいても平方完成は有用である.

      ここから学ぶべきことは,定数項のやり取りによって式を無理やりわかりやすい形,知っている形,扱いやすい形に変形できるということだ.

    \subsection{部分分数分解}
      先程は和による値の調整だったが,積で調整する場合を考えよう.
      \[
        \frac{1}{(x+a)(x+b)}
      \]
      のような値は,シグマを用いた計算をする上で扱いづらい.
      そうでなくとも分母の次数はなるべく簡単な数字であってほしい.
      そのような要望から,以下のような変形を考える.
      \[
        \frac{1}{(x+a)(x+b)}=\frac{\alpha}{x+a}+\frac{\beta}{x+b}
      \]
      $\alpha,\beta$について求めると,
      \begin{align*}
      \frac{1}{(x+a)(x+b)}
        &= \frac{\alpha}{x+a}+\frac{\beta}{x+b}\\
        &= \frac{(\alpha+\beta)x+\alpha b+\beta a}
        {(x+a)(x+b)}
      \end{align*}
      \[
      \therefore\quad
      \begin{cases}
      \alpha+\beta=0,\\
      \alpha b+\beta a=1
      \end{cases}
      \]
      この連立方程式を解けば1次の分数の和として表すことができるのだ.
      3次の場合や$x$に係数がある場合でも
      \[
        \frac{1}{(px+a)(qx+b)(rx+c)}=\frac{\alpha}{px+a}+\frac{\beta}{qx+b}+\frac{\gamma}{rx+c}
      \]
      とすることで解を得ることができる.

